% arara: lualatex: { synctex: on }
% arara: biber
% arara: lualatex: { synctex: on }
% arara: lualatex: { synctex: on }
\documentclass[titlepage,a4paper,12pt]{article}

% ==============================================================================
% 1. Detección de motor TeX y soporte de fuentes nativas (LuaLaTeX / XeTeX)
% ==============================================================================
\usepackage{iftex}
\usepackage{fontspec}
\usepackage{metalogo}       % Proporciona los logos de \LuaTeX, \XeTeX, etc.

% Configuración de fuentes del documento
\defaultfontfeatures{Ligatures=TeX}
\setmainfont[mode=harf,script=latin]{Tex Gyre Termes}
\setsansfont[mode=harf,script=latin]{Arial}
\setmonofont[mode=harf,Color={0019D4},script=latin]{Courier New}

% ==============================================================================
% 2. Geometría y estructura de página
% ==============================================================================
\usepackage[margin=1in]{geometry}
\usepackage{fancyhdr}
\usepackage{setspace}
\setstretch{1.5}            % Interlineado 1.5 limpio mediante setspace

% ==============================================================================
% 3. Gráficos, colores y tablas
% ==============================================================================
\usepackage{xcolor}         % Siempre antes de tikz para evitar option clashes
\usepackage{graphicx}       % Carga única (se eliminó el duplicado anterior)
\usepackage{tikz}
\usetikzlibrary{shadows}
\usepackage{pgfgantt}
\usepackage{float}
\usepackage{wrapfig}
\usepackage{subcaption}
\usepackage{tabularx}
\usepackage{tabto}

% ==============================================================================
% 4. Utilidades generales, listas y formato de títulos
% ==============================================================================
\usepackage{import}
\usepackage{paralist}
\usepackage{widows-and-orphans}
\usepackage{etoolbox}
\usepackage{keyval}
\usepackage{ifthen}
\usepackage{titlesec}
\titleformat{\section}
  {\normalfont\sffamily\large\bfseries}
  {\thesection}{1em}{}

% Paquete para citas de textos bíblicos
\usepackage{scripture}

% ==============================================================================
% 5. Idiomas, tipografía y comillas
% ==============================================================================
\usepackage{polyglossia}
\setmainlanguage{spanish}
\setotherlanguage[variant=ancient]{greek}
\setotherlanguage{hebrew}
\setotherlanguage{english}
\SetLanguageKeys{spanish}{indentfirst=false}

% Fuentes para idiomas bíblicos/antiguos
\newfontfamily\greekfont[mode=harf,Script=Greek]{SBL_grk.ttf}
\newfontfamily\hebrewfont[mode=harf,Script=Hebrew]{SBL_Hbrw.ttf}

% Tipografía mejorada y comillas contextuales (csquotes debe ir DESPUÉS de polyglossia)
\usepackage{microtype}
\usepackage[autostyle=true]{csquotes}

% ==============================================================================
% 6. Bibliografía y apéndices (deben ir DESPUÉS de csquotes y ANTES de hyperref)
% ==============================================================================
\usepackage[notes,backend=biber,isbn=false,ibidtracker=constrict,idemtracker=constrict,useibid=true]{biblatex-chicago}
\renewcommand{\bibsetup}{\singlespacing}
\renewcommand{\bibitemsep}{1\baselineskip}
\renewcommand{\bibhang}{0.5in}
\addbibresource{template.bib}

\usepackage[toc,page]{appendix}
\renewcommand{\appendixpagename}{Apéndices}
\renewcommand{\appendixtocname}{Apéndices}

% Macros auxiliares para la tabla de contenidos
\newcommand{\stoptocwriting}{%
  \addtocontents{toc}{\protect\setcounter{tocdepth}{-5}}}
\newcommand{\resumetocwriting}{%
  \addtocontents{toc}{\protect\setcounter{tocdepth}{\arabic{tocdepth}}}}

% ==============================================================================
% 7. Hipervínculos, referencias y paquetes dependientes de hyperref
% (hyperref debe ser de los últimos, seguido estrictamente por bookmark, cleveref y pageslts)
% ==============================================================================
\usepackage{url}
\usepackage[hidelinks,colorlinks=false,linkcolor=black,urlcolor=black,bookmarksopen=true,linktoc=all]{hyperref}
\usepackage{bookmark}
\usepackage{cleveref}
\usepackage{pageslts}       % Requiere cargarse después de hyperref para el conteo de páginas

% ==============================================================================
% 8. Configuración de encabezados, pie de página y metadatos
% ==============================================================================
\pagestyle{fancy}% Encabezado y pie de página
\fancyhf{}
\fancyhead[L]{\textsf{\title}}
\fancyhead[R]{\textsf{\nombremateria}}
\fancyfoot[C]{\textsf{\thepage\ de~\pageref*{VeryLastPage}}}
\setlength{\headheight}{14.5pt}

\renewcommand{\footnotesize}{\fontsize{10pt}{12pt}\selectfont}

\def\title{Titulo del trabajo}
\def\author{Nombre del autor}
\def\instructor{Nombre del instructor}
\def\fechadeentrega{Fecha de entrega}
\def\nombremateria{Nombre de Materia}
\def\codigomateria{Codigo de materia}

% ==============================================================================
% Cuerpo del Documento
% ==============================================================================
\begin{document}
\pagenumbering{roman}
\begin{titlepage} % Suppresses displaying the page number on the title page and the subsequent page counts as page 1
	\newcommand{\HRule}{\rule{\linewidth}{0.5mm}} % Defines a new command for horizontal lines, change thickness here

	\center% Centre everything on the page

	%------------------------------------------------
	%	Logo
	%------------------------------------------------
	\includegraphics[width=6cm]{./FTIBA-logo.png}\\

	%------------------------------------------------
	%	Title
	%------------------------------------------------

	\HRule\\[0.4cm]

	{\huge\bfseries \title}\\[0.4cm] % Title of your document

	\HRule\\[1.5cm]
	%------------------------------------------------
	%	Author(s)
	%------------------------------------------------

	\begin{minipage}{0.4\textwidth}
		\begin{flushleft}
			\large
			\textit{Author}\\
			\textsc{\author} % Your name
		\end{flushleft}
	\end{minipage}


	% If you don't want a supervisor, uncomment the two lines below and comment the code above
	%{\large
	%\textit{Autor}}\\
	%\Large
	%\textsc{\author} % Your name
	\vfill
	%------------------------------------------------
	%	Headings
	%------------------------------------------------
	\textsc{\Large \codigomateria}\\[0.5cm] % Major heading such as course name

	\textsc{\large \nombremateria}\\[0.5cm] % Minor heading such as course title
	\vfill


	\large
	\textit{Instructor}\\
	\textsc{\instructor} % Supervisor's name
	%------------------------------------------------
	%	Date
	%------------------------------------------------

	\vfill\vfill\vfill % Position the date 3/4 down the remaining page

	{\large \fechadeentrega} % Date, change the \today to a set date if you want to be precise


	%\vfill\vfill
	%\includegraphics[width=0.2\textwidth]{placeholder.jpg}\\[1cm] % Include a department/university logo - this will require the graphicx package

	%----------------------------------------------------------------------------------------

	\vfill % Push the date up 1/4 of the remaining page

\end{titlepage}

\thispagestyle{empty}
\tableofcontents
\clearpage

\pagenumbering{arabic}
\setlength{\parskip}{2mm}
\section{Introducción}
\import{./}{primera-seccion.tex}
\import{./}{segunda-seccion.tex}
\import{./}{conclusion.tex}
\clearpage
\printbibliography[heading=bibintoc, title={Trabajos Citados}]
\appendixpage
\begin{appendix}
\end{appendix}
\end{document}
